% Timaeus --- headnote
% timaeus, 45 BC, Rome

\textit{Cicero's \textit{Timaeus} is a Latin rendering of the central
cosmological discourse of Plato's dialogue of that name --- the account
of how the world was made --- undertaken as part of his campaign to give
Rome a philosophical literature of its own. It survives as a fragment,
and its exact date and purpose are uncertain; it may have been meant as
raw material for a dialogue of Cicero's own, or as a companion to other
works of these years. It opens, in any case, with a setting of Cicero's
own invention: a meeting at which he falls into philosophical
conversation with Publius Nigidius Figulus, the most learned Roman of the
age after Varro and a reviver of Pythagorean studies, and with Cratippus
the Peripatetic, whom Cicero counted the first philosopher of his time.}

\textit{From the third section the text becomes Plato: the discourse of
Timaeus of Locri on the origin of the universe. There is that which
always is and never comes to be, and that which is always coming to be
and never is; the world, being visible and tangible, belongs to the
second kind, and so must have had a maker and a beginning. That maker,
being good and free of envy, framed the world upon the model of the
eternal pattern, made it a single living creature endowed with soul and
mind, and shaped its body out of fire and earth bound fast by the
proportions of air and water. Cicero follows Plato through the
mathematics of the world-soul, the harmonic intervals by which it is
divided, and the circles of the same and the different that carry the
heavens --- the most demanding stretch of the work, where Latin is made
to do what it had never done before.}

\textit{The chief interest of the piece is precisely that strain. Cicero
is inventing a philosophical Latin as he goes, coining words and
stretching old ones to carry Greek metaphysics, and he says as much in
passing --- that one must be bold, since these things are here first
being named anew. Plato's \textit{Timaeus}, in the later Latin version of
Calcidius, would become the one dialogue the Latin Middle Ages knew at
first hand and the seedbed of medieval cosmology; Cicero's earlier and
partial rendering is the first attempt to bring that vision into Latin at
all, and a rare chance to watch a great writer of prose wrestling another
language's deepest abstractions into his own.}
